%********************************************************************
% Appendix
%*******************************************************
% If problems with the headers: get headings in appendix etc. right
%\markboth{\spacedlowsmallcaps{Appendix}}{\spacedlowsmallcaps{Appendix}}
\chapter[Expérience annexe: période d'attention]{Séquence d'événements ou texture sonore: l'influence de la période d'attention.}
\label{app:xp_texture}

\section{Objectif de l'expérience}

Nous présentons ici les résultats d'une étude sur la perception des textures sonores, étude menée dans le cadre de cette thèse, mais déconnectée du sujet principal. 

Comme nous l'avons vu (\cf~Section~\ref{sec:ch3_textureDef}), la texture est un objet composite, dont les éléments constitutifs cessent d'être perçus de manière distincte, dès lors qu'ils occurrent suivant un pattern dont les caractéristiques physiques restent stables au cours du temps. La perception de ce pattern nécessite une certaine période d'attention. Nous proposons ici un protocole expérimental permettant d'étudier ces périodes d'attention.


En considérant comme stimuli une mixture d'événements du même type, nous faisons l'hypothèse qu'à partir du moment où le cerveau parvient à isoler un événement de cette mixture, il ne perçoit plus la mixture comme une texture, mais comme une succession d'événements. Inversement, s'il ne parvient pas à distinguer un événement isolé, alors la mixture est perçue comme une texture.

Nous appliquons le modèle de scène sonore proposé (\cf~Section~\ref{sec:ch4_modelForm}) afin de simuler des textures à partir de séquences d'événements dont nous contrôlons l'espacement \emph{inter-onsets} moyen. Il s'agit de faire varier cet espacement, afin d’identifier le seuil à partir duquel la séquence d'événements cesse d'être perçue comme une texture. 

À ce titre, le protocole proposé s'inscrit complètement dans le cadre des expériences perceptives portant sur la détection du signal.

\section{Banque de données}


Chaque stimuli est composé d'un son cible, suivi d'une séquence d'événements enchevêtrés. Le son cible est le même pour tous les stimuli, et tous les sujets. Il a été choisi par les expérimentateurs, et voulu ni trop identifiable, ni tout à fait dénué de caractéristiques saillantes. 

Les événements sont tous des sons isolés ayant une durée de $1$ seconde. Les séquences durent 6 secondes. Chacune donne à entendre des scènes de trafic. Ces scènes sont simulées en agglomérant des sons de voiture (isolés). La simulation est contrôlée par un paramètre réglant l'espacement temporel inter-\emph{onsets} moyen entre les événements. Cinq valeurs d'espacement sont considérées: $0.1$, $0.3$, $0.5$, $0.7$ et $0.9$ secondes (\cf~Figure~\ref{fig:xptexturea}, et~\ref{fig:xptextureb}). 

Pour chaque espacement, nous simulons 20 séquences de trafic. Les sujets doivent donc écouter 100 stimuli. La moitié d'entre eux sont des pièges (\emph{catch trial}). Ils ne contiennent pas le son cible.

\section{Planification expérimentale}

{\setlength{\parindent}{0cm}\textbf{Procédure}} \\ 

L'expérience est une épreuve d'évaluation de type oui/non. Chaque sujet évalue l'ensemble des 100 stimuli. Pour chacun il indique si oui ou non il entend le son cible.

Les scènes sont présentées dans un ordre aléatoire.  \\

{\setlength{\parindent}{0cm}\textbf{Dispositif expérimental}} \\

Tous les sujets passent l'expérience sur des machines identiques. L'audio est diffusé en monophonique, par le biais de casques audio semi-ouverts \emph{Beyer-Dynamic DT 990 Pro}. 

Le niveau sonore de sortie est le même pour tous les sujets. Il a été préalablement fixé par les expérimentateurs afin de correspondre à un niveau d'écoute confortable.

Les sujets passent l'expérience individuellement, dans un environnement acoustique calme. Un expérimentateur est présent durant la totalité de l'expérience, afin de contrôler le bon déroulement de cette dernière, et de répondre aux éventuelles questions des sujets. \\

{\setlength{\parindent}{0cm}\textbf{Participants}} \\

13 sujets (5 femmes), tous étudiants à l'École Centrale de Nantes, ont réalisé l'expérience.
 
\section{Méthodologie et outils statistiques}
\label{app:xptextureMethodo}

\begin{table}[t]
\centering
\begin{tabular}{c| c c}  
                         & réponse: oui & réponse: non  \\      
\hline
Présence son cible: oui  & $vp$           & $fr$ \\
Présence son cible: non  & $fp$           & $vr$ \\
\hline
\end{tabular}
\vspace{0.5mm}
\caption[Théorie de la détection du signal.]{Théorie de la détection du signal.}
\label{tab:sdtHitFP}
\end{table}

Pour chaque condition expérimentale (espacement inter-\emph{onsets}), les performances des sujets sont évaluées à l'aide de la mesure de sensitivité $d'$.

\begin{equation}
d'=z(H)-z(F)
\end{equation}

\begin{equation}
H=\dfrac{vp}{vp+fr} \quad \textrm{,} \quad F=\dfrac{fp}{fp+vr}
\end{equation}

avec $H$ le taux de réponses correctes, $F$ le taux de faux positifs, $vp$ le nombre de vrais positifs, $fr$ le nombre de fausses réjections, $fp$ le nombre de fausses alarmes, $vr$ le nombre de vraies réjections (\cf~Tableau~\ref{tab:sdtHitFP}), et $z(x)$ l'opérateur de transformation $z$ Gaussien permettant d'exprimer la quantité $x$ en nombre d'écarts types par rapport à la moyenne.
L'analyse se déroule en deux temps.

\begin{itemize}
\item Pour apprécier l'existence de différences au niveau des $d'$ entre les conditions expérimentales, nous considérons une analyse de variance à mesures répétées comportant une variable indépendante à 5 niveaux: l'espacement inter-\emph{onsets}. La sphéricité est évaluée à l'aide du test de Mauchly. Si cette dernière n'est pas vérifiée, la correction de Greenhouse-Geisser est appliquée, la valeur $p$ est alors notée $p_{gg}$. Le seuil de significativité est fixé à $\alpha=0.05$;
\item Afin d'établir un seuil théorique d'espacement en deçà duquel les sujets n'arrivent plus à distinguer le son cible, nous calculons une régression linaire des $d'$ en fonction de l'espacement, et considérons la valeur de l'espacement pour $d'=1$. Ce seuil de $d'=1$ est couramment utilisé dans la théorie de la détection du signal \citep{macmillan2004detection}.
\end{itemize}

\section{Résultats}

\begin{figure}[t]
        \myfloatalign
        \subfloat[]
        {\includegraphics[width=.5\linewidth]{gfx/app_texture/xpTexture0_1}\label{fig:xptexturea}}
        \subfloat[]
        {\includegraphics[width=.5\linewidth]{gfx/app_texture/xpTexture0_2}\label{fig:xptextureb}} \par
        \subfloat[]
        {\includegraphics[width=.5\linewidth]{gfx/app_texture/xpTexture_1}\label{fig:xptexturec}}
        \subfloat[]
        {\includegraphics[width=.5\linewidth]{gfx/app_texture/xpTexture_2}\label{fig:xptextured}}
        \caption[Événement ou texture sonore: influence de la période d'attention.]{Événement ou texture sonore: influence de la période d'attention. (a) stimulus ayant un espacement inter-\emph{onsets} élevé; (b) stimulus ayant un espacement inter-\emph{onsets} faible; (c) évolution des taux de réponses correctes $H$ et de faux positif $F$ en fonction de l'espacement; (d) régression linaire de $d'$ en fonction de l'espacement.}\label{fig:xptexture}
\end{figure}

L'ANOVA à mesures répétées montre un effet significatif de l'espacement sur la sensitivité $d'$ ($F[4,48]=24$, $p_{gg}<0.01$). L'observation des évolutions des taux $H$ et $F$ montre que les réponses des sujets ne sont pas biaisées: le taux de réponses correctes $H$ (\ie, la capacité à distinguer le son cible) continue d'augmenter avec l'espacement, tandis que le taux de faux positif $F$ reste à peu près constant à partir d'une valeur d'espacement de 0.3 seconde (\cf~Figure~\ref{fig:xptexturec}).

La régression linaire de $d'$ en fonction de l'espacement (\cf~Figure~\ref{fig:xptextured}) nous permet d'obtenir une limite théorique ($d'=1$) à partir de laquelle la mixture cesse d'être perçue comme une texture. Pour des sons isolés d'une seconde, cet espacement limite est de 0.46 seconde, soit la moitié de la durée des événements utilisés.

 


